\documentclass[11pt,a4paper]{article}
\usepackage[left=3cm,right=3cm,top=3cm,bottom=3cm]{geometry}
\usepackage{lmodern}
\usepackage{setspace}
\usepackage{parskip}
\usepackage{hyperref}
\usepackage{amsmath, amssymb}
\usepackage{graphicx}
\usepackage{biblatex}
\addbibresource{references1.bib}

\title{Research Interests in Industrial Organization and Business Modeling}
%  \\
% \large Their Relevance to the Chair of International Management
\author{%
Zhongli Wang \\\\%
% \href{mailto:nkwangzhongli@gmail.com}{nkwangzhongli@gmail.com} \\\\%
% Bielefeld University%
}
% \date{\today}


\begin{document}

\maketitle

% \begin{abstract}
% This short essay outlines my core research interests in Industrial Organization (IO) and economic modeling, and illustrates how these themes align with the strategic and methodological focus of the Chair of International Management. I reflect on past and ongoing research, methodological approaches, and future directions for collaboration.
% \end{abstract}

\section*{Introduction}
Industrial Organization (IO) examines how firms interact in markets, how market structure affects economic outcomes, and how strategic behavior evolves under competition. My academic background and doctoral research are rooted in these themes, with a particular focus on agent-based simulations and dynamic firm strategies.

\section*{Research Background and Interests}
In my dissertation, titled \textit{Strategic Behavior and Economic Dynamics in ICT Markets: An Agent-Based Framework}, I explore how firms adjust pricing, investment, and capacity over time in response to technological change and consumer heterogeneity. This work integrates IO theory with computational modeling to analyze market outcomes under various policy and strategic scenarios.

My broader interests include:
\begin{itemize}
  \item Firm dynamics under imperfect competition
  \item Network effects and switching costs
  \item Heterogeneous agents and market evolution
  \item Policy analysis using simulation-based models
  \item Applied machine learning in digital and financial market
  % \item Digital economics 
  % \item Financial market
\end{itemize}

\section*{Methodological Approach}
I combine theoretical and computational tools to investigate market outcomes. These include:
\begin{itemize}
  \item Agent-Based Modeling (ABM) for firm and consumer interactions
  \item Simulation-based estimation techniques (e.g., Simulated Method of Moments)
  \item Network-based modeling of consumer behavior
\end{itemize}

Such methods allow flexible modeling of heterogeneity, strategic learning, and non-linear features that are often crucial in real-world strategic environments.

\section*{Relevance to the Chair of International Management}
The Chair's emphasis on strategic management, international research orientation, and methodological rigor suits my research profile. In particular, my experience in simulation methods and firm behavior under uncertainty aligns with ongoing work at the Chair.

I am especially interested in potential collaboration on:
\begin{itemize}
  \item Strategic decision-making in global firms
  \item Firm behavior and market evolution in the financial market
  \item Competitive dynamics and technological adaptation
  \item Methodological innovation in simulation and empirical validation
\end{itemize}

\section*{Proposed Research Questions}
\textbf{RQ 1:} How do multinational enterprises (MNEs) strategically choose and adapt the locations of their operations within increasingly fragmented and volatile global value chains?

Recent geopolitical tensions, trade policy shifts, and post-pandemic disruptions have brought location strategy to the forefront of managerial and policy discussions. Issues such as reshoring, nearshoring, and supply chain resilience now dominate both corporate and public policy agendas (e.g., Economist, 2023; JIBS, 2022). Despite the centrality of these questions, there remains a lack of dynamic, firm-level models that capture how MNEs respond over time to economic, institutional, and regulatory pressures in both domestic and international contexts. Building on my background in computational modeling and industrial organization, I propose to develop agent-based simulation frameworks that examine how firms explore, commit to, and reconfigure location choices under policy uncertainty, labor market shifts, and infrastructure variation. This line of research aims to bridge theoretical insights from strategic management—particularly in dynamic capabilities and organizational renewal—with practical concerns faced by MNEs navigating evolving global environments. It would also contribute to a better understanding of how macro-level policy interventions (e.g., trade agreements, subsidies, decarbonization policies) shape micro-level strategic behavior and long-term performance outcomes.

\textbf{RQ 2:} How do firms in emerging technology sectors—such as robotics and large language models—navigate strategic renewal under uncertainty and performance pressure?

Building on Keil et al. (2008), I propose to develop an agent-based simulation model that examines how firms select governance modes for external innovation—corporate venture capital (CVC), alliances, joint ventures, or acquisitions—based on inter-organizational trust, technological relatedness, and environmental turbulence. At the same time, drawing on insights from Posen et al. (2017), I plan to integrate a behavioral logic of problemistic search, where firms only adapt or initiate governance changes when performance falls below aspiration levels. The simulation will capture how firms experiment, imitate, or renew their strategies in response to perceived underperformance, while facing evolving market dynamics and innovation pressure. Preliminary insights suggest that firms relying solely on static governance routines may miss windows of opportunity or overcommit to misaligned innovation modes. By combining the behavioral theory of the firm with strategy process research on governance and external venturing, this project contributes to a deeper understanding of how firms reconfigure their innovation strategies through adaptive governance under uncertainty.

\textbf{RQ 3:} Why do some technology firms choose not to internationalize? 

A third line of inquiry focuses on the strategic decision-making processes behind internationalization — or the lack thereof — among technology firms. While much of the literature emphasizes how multinational enterprises (MNEs) compete globally, relatively less attention is paid to why some high-potential firms deliberately choose to remain domestic, especially in highly scalable sectors like ICT and AI. Drawing on behavioral strategy and international business theory, this project aims to explore how firms interpret market signals and make internationalization decisions under uncertainty, bounded rationality, and institutional complexity. I plan to develop an agent-based simulation in which heterogeneous firms decide whether and when to enter foreign markets, based on noisy signals about user demand, regulatory environments, and competitor activity. These simulations will allow us to examine how cognitive limitations, organizational routines, or risk-aversion influence strategic inertia. Preliminary results suggest that even capable firms may rationally delay or avoid international expansion due to misperceptions, lack of absorptive capacity, or uncertainty amplification. This research contributes to the strategic management literature by explaining non-entry and delayed internationalization not as failures, but as adaptive responses to complex global environments — with implications for policymakers and support programs targeting SME internationalization.

\textbf{RQ 4:} How do investor expectations and capital market dynamics influence strategic behavior and survival of robotics firms in emerging industries?

A fourth research direction explores the relationship between investor expectations and firm strategy in the context of emerging technologies, focusing specifically on the rise and retreat of investment in the robotics industry. In recent years, numerous robotics firms in China and beyond have experienced rapid early-stage growth supported by both state incentives and speculative private capital. However, many of these firms faced abrupt funding contractions as investor enthusiasm declined, resulting in strategic pivots, downsizing, or exit. This pattern raises important questions about how firms formulate and adjust their strategies under volatile capital conditions, and how investor–firm interactions shape competitive outcomes in nascent industries. Using an agent-based simulation framework, I plan to model the co-evolution of firm decision-making and investor behavior, capturing how hype cycles, performance signals, and policy support influence strategic renewal, overexpansion, and market consolidation. Preliminary results suggest that short-term capital can distort firm incentives, leading to premature scaling and industry fragility, whereas firms supported by long-horizon investors or adaptive resource allocation strategies are more likely to survive and lead in the long term. This research contributes to the strategic management literature on resource dependence, strategic flexibility, and innovation ecosystems, while offering insights into how policy and capital structures affect the sustainability of industrial emergence.

\textbf{RQ 5:} How do multinational firms in the AI software sector compete under data and algorithmic regulation, considering network effects and heterogeneous user segments?

As firms like OpenAI, DeepSeek, and Mistral compete across regions, they must navigate powerful network effects, heterogeneous user preferences, and an increasingly complex landscape of AI regulation. Policies such as the EU AI Act, China’s algorithm oversight rules, and U.S. export controls are likely to shape both innovation and market access, raising critical questions about the international management of digital platforms. This project aims to explore how MNEs adapt their strategies—regarding pricing, localization, and compliance—in response to regulatory fragmentation and shifting consumer segments. Using a combination of agent-based modeling, user-level behavioral simulation, and machine learning tools, I plan to investigate how firms develop competitive advantages in data-sensitive, regulation-driven global markets. This research would contribute to strategic management literature on technological competition, global value chain coordination, and organizational adaptation under regulatory constraints, while offering practical insights into how LLM firms can align strategic decision-making with both market dynamics and public policy.

\section*{Conclusion and Outlook}
My goal is to further develop a rigorous research portfolio in applied IO, combining theory, computation, and empirical relevance. I see strong opportunities for synergy with the Chair's focus areas, and I am eager to contribute to an ambitious and collaborative research environment.


\printbibliography
\end{document}

% \vspace{2em}
% \noindent\textbf{Keywords:} Industrial Organization, Agent-Based Modeling, Firm Strategy, Economic Dynamics, Simulation Methods